% Source: PDF pages 36--37.
\subsection{Expected number of leaves using Lagrange inversion}
Suppose $F(u,z)$ is a bivariate generating function, with parameter $\chi$. Recall that
\[
\mathbb E_n[\chi]
=\frac{[z^n]\left.\dfrac{\partial}{\partial u}F(u,z)\right|_{u=1}}
{[z^n]F(1,z)}.
\]
Coefficient extraction and differentiation commute:
\[
[z^n]\left.\frac{\partial}{\partial u}F(u,z)\right|_{u=1}
=\left.\frac{\partial}{\partial u}[z^n]F(u,z)\right|_{u=1}.
\]
% Source page 36 writes F(u,z) rather than F(1,z) in the denominator
% of the expectation formula; evaluation at u=1 is required.
\begin{example}
The bivariate generating function for nonempty binary trees, enumerated by size and number of leaves, satisfies
\[
N(u,z)=uz+2zN(u,z)+zN(u,z)^2.
\]
Here the $uz$ term represents a single leaf. Thus
\[
z=\frac{N(u,z)}{G_u(N(u,z))},\qquad G_u(t)=u+2t+t^2.
\]
Using Lagrange inversion,
\begin{align*}
[z^n]\left.\frac{\partial}{\partial u}N(u,z)\right|_{u=1}
&=\left.\frac{\partial}{\partial u}[z^n]N(u,z)\right|_{u=1}\\
&=\left.\frac{\partial}{\partial u}
\left(\frac1n[t^{n-1}](u+2t+t^2)^n\right)\right|_{u=1}\\
&=\frac1n[t^{n-1}]
\left.\frac{\partial}{\partial u}(u+2t+t^2)^n\right|_{u=1}\\
&=\frac1n[t^{n-1}]\left.n(u+2t+t^2)^{n-1}\right|_{u=1}\\
&=[t^{n-1}](1+2t+t^2)^{n-1}\\
&=[t^{n-1}](1+t)^{2(n-1)}
=\binom{2n-2}{n-1}.
\end{align*}
Consequently, for $n\ge1$,
\begin{align*}
\mathbb E_n[\chi]
&=\frac{\binom{2n-2}{n-1}}{\frac1{n+1}\binom{2n}{n}}\\
&=(n+1)\frac{(2n-2)!}{(2n)!}
\frac{n!}{(n-1)!}\frac{n!}{(n-1)!}\\
&=\frac{n(n+1)}{2(2n-1)}\sim\frac n4.
\end{align*}
Thus asymptotically one quarter of the vertices in a uniformly chosen binary tree of size $n$ are leaves, on average.
\end{example}
